% !TeX root = part5.tex
\documentclass[../final.tex]{subfiles}

\begin{document}

\ifSubfilesClassLoaded{}{\addtocontents{toc}{\protect\newpage}}
\practice{5}{01.10.2026}{Задачи по физике детекторов}

\task[1]{Пробег B-мезонов и асимметрия коллайдера}
\label{task:b_decay_length_sem5}

\condition
При энергии в системе центра масс $E_{\text{с.ц.м.}}=10.58\,\text{ГэВ}$
образуется пара $B\overline B$. Масса мезона $m_Bc^2=5.28\,\text{ГэВ}$,
собственная длина распада $c\tau\simeq491.1\,\text{мкм}$ (для $B^\pm$).
Оценить средний пробег мезона в системе центра масс и в лабораторной системе
для KEKB с энергиями пучков $8$ и $3.5\,\text{ГэВ}$.
Объяснить, зачем нужна асимметрия пучков.

\solution
\subsection*{Движение мезонов в системе центра масс}

Двум мезонам одинаковой массы достаются равные энергии:
\begin{equation*}
    E_B'=\frac{E_{\text{с.ц.м.}}}{2}=5.29\,\text{ГэВ}.
\end{equation*}
Штрихом обозначаем величины в системе центра масс. Импульс равен
\begin{equation*}
    p_B'c=\sqrt{(E_B')^2-m_B^2c^4}
    =\sqrt{5.29^2-5.28^2}\simeq0.325\,\text{ГэВ}
    \simeq0.33\,\text{ГэВ}.
\end{equation*}
Поэтому
\begin{equation*}
    \beta_B'=\frac{p_B'c}{E_B'}\simeq0.061,
    \qquad
    \gamma_B'=\frac{E_B'}{m_Bc^2}\simeq1.002.
\end{equation*}
Из $p=\gamma mv$ также следует
\begin{equation*}
    \beta_B'=\frac{p_B'c}{\gamma_B'm_Bc^2}
    \simeq\frac{p_B'c}{m_Bc^2},
\end{equation*}
где последнее приближение допустимо из-за $\gamma_B'\simeq1$.
Средний пробег до распада:
\begin{equation*}
    \ell_B'=\beta_B'\gamma_B'c\tau
    =\frac{p_B'c}{m_Bc^2}c\tau
    \simeq30\,\text{мкм}.
\end{equation*}
При симметричном коллайдере лабораторная система совпадает с системой
центра масс, поэтому вершины распада находятся очень близко к точке столкновения.

\begin{rbox}[left=16pt,right=16pt,top=12pt,bottom=10pt]{[]}
    \centering
    \begin{tikzpicture}[>=Latex,every node/.style={text=rtext}]
        \fill[rtext] (0,0) circle(2pt);
        \draw[rc3,thick,->] (0,0)--(2,0.8) node[right] {$B$};
        \draw[p3,thick,->] (0,0)--(2,-0.8) node[right] {$\overline B$};
        \fill[rc3] (1.7,0.68) circle(2pt);
        \fill[p3] (1.4,-0.56) circle(2pt);
        \draw[rc3,->] (1.7,0.68)--(2.8,1.5);
        \draw[rc3,->] (1.7,0.68)--(3.0,0.9);
        \draw[p3,->] (1.4,-0.56)--(2.8,-0.7);
        \draw[p3,->] (1.4,-0.56)--(2.4,-1.5);
        \draw[<->,rtextmuted] (1.4,-1.85)--(1.7,-1.85);
        \node[below] at (1.55,-1.85) {$\Delta z$};
        \node[left,align=right] at (0,0) {точка\\столкновения};
    \end{tikzpicture}
    \captionsetup{font=footnotesize,labelfont=bf,skip=5pt}
    \captionof{figure}{Схема разнесённых вершин распада в лабораторной системе.
    Измерение $\Delta z$ позволяет оценить разность времён распада. Масштаб условный.}
    \label{fig:b_vertices_sem5}
\end{rbox}

\subsection*{Движение системы центра масс в лабораторной системе}

Для встречных ультрарелятивистских пучков
\begin{equation*}
    E_{\text{лаб}}=E_{e^-}+E_{e^+}=11.5\,\text{ГэВ},
    \qquad
    P_{\text{лаб}}c=E_{e^-}-E_{e^+}=4.5\,\text{ГэВ}.
\end{equation*}
Параметры преобразования из системы центра масс в лабораторную:
\begin{equation*}
    \gamma_{\text{ц.м.}}=\frac{11.5}{10.58}\simeq1.09,
    \qquad
    \beta_{\text{ц.м.}}=\frac{4.5}{11.5}\simeq0.39.
\end{equation*}
Пренебрегая малой скоростью мезона в системе центра масс,
оцениваем его средний продольный пробег:
\begin{equation*}
    \ell_{\text{лаб}}\simeq
    \beta_{\text{ц.м.}}\gamma_{\text{ц.м.}}c\tau
    \simeq0.39\cdot1.09\cdot491.1\,\text{мкм}
    \simeq210\,\text{мкм}.
\end{equation*}

\answer
В системе центра масс средний пробег составляет около
\rresult{$30\,\text{мкм}$}, в лабораторной системе --- порядка
\rresult{$0.21\,\text{мм}$}. Асимметрия энергий пучков увеличивает
пространственное разделение вершин и облегчает измерение разности времён
распада:
\begin{equation*}
    \Delta t\simeq\frac{\Delta z}
    {\beta_{\text{ц.м.}}\gamma_{\text{ц.м.}}c}.
\end{equation*}

\task[2]{Пороги черенковского детектора}
\label{task:cherenkov_thresholds_sem5}

\condition
Найти пороговые импульсы пионов и каонов в аэрогеле ($n=1.01$)
и кварце ($n=1.54$). Принять
\begin{equation*}
    m_\pi c^2\simeq140\,\text{МэВ},\qquad
    m_Kc^2\simeq494\,\text{МэВ}.
\end{equation*}

\solution
Черенковское излучение возникает при $\beta>1/n$.
На пороге
\begin{equation*}
    \beta_{\text{пор}}=\frac1n,\qquad
    \gamma_{\text{пор}}=\frac1{\sqrt{1-1/n^2}}.
\end{equation*}
Используя $pc=\gamma\beta mc^2$, получаем
\begin{equation*}
    (pc)_{\text{пор}}=\frac{mc^2}{\sqrt{n^2-1}}.
\end{equation*}

\subsection*{Аэрогель}

При $n=1.01$
\begin{equation*}
    \beta_{\text{пор}}\simeq0.990,
    \qquad \gamma_{\text{пор}}\simeq7.12.
\end{equation*}
Пороговые импульсы:
\begin{equation*}
    p_{\pi,\text{пор}}\simeq0.99\,\text{ГэВ}/c,
    \qquad
    p_{K,\text{пор}}\simeq3.48\,\text{ГэВ}/c
    \simeq3.5\,\text{ГэВ}/c.
\end{equation*}

\subsection*{Кварц}

При $n=1.54$
\begin{equation*}
    \beta_{\text{пор}}\simeq0.649,
    \qquad\gamma_{\text{пор}}\simeq1.315.
\end{equation*}
Поэтому
\begin{equation*}
    p_{\pi,\text{пор}}\simeq0.120\,\text{ГэВ}/c,
    \qquad
    p_{K,\text{пор}}\simeq0.422\,\text{ГэВ}/c.
\end{equation*}

\answer
\begin{center}
    \begin{tabular}{lccc}
        \hline
        \rconcept{Среда} & $n$ & $p_{\pi,\text{пор}},\ \text{ГэВ}/c$
        & $p_{K,\text{пор}},\ \text{ГэВ}/c$\\
        \hline
        Аэрогель & 1.01 & 0.99 & 3.48\\
        Кварц & 1.54 & 0.120 & 0.422\\
        \hline
    \end{tabular}
\end{center}
В аэрогеле при импульсах между $0.99$ и $3.48\,\text{ГэВ}/c$
пион излучает, а каон ещё не излучает. Это позволяет различать их
при одном и том же измеренном импульсе.

\task[3]{Восстановление частиц по фотонам в калориметре}
\label{task:calorimeter_sem5}

\condition
В калориметре зарегистрированы четыре вспышки. Считаем их фотонами.
Их энергии:
\begin{equation*}
    E_1=1\,\text{ГэВ},\quad E_2=1.5\,\text{ГэВ},\quad
    E_3=1.7\,\text{ГэВ},\quad E_4=0.5\,\text{ГэВ}.
\end{equation*}
Углы между направлениями:
\begin{align*}
    \theta_{12}&=6.3^\circ,& \theta_{13}&=12.7^\circ,&
    \theta_{14}&=160^\circ,\\
    \theta_{23}&=23.5^\circ,& \theta_{24}&=85.0^\circ,&
    \theta_{34}&=34.6^\circ.
\end{align*}
Найти инвариантные массы пар и определить возможные родительские частицы.

\solution
\rassumption{Используем единицы $c=1$}.
Для любых двух частиц
\begin{equation*}
    m_{ij}^2=m_i^2+m_j^2+2E_iE_j
    -2p_ip_j\cos\theta_{ij}.
\end{equation*}
Для фотонов $m_i=m_j=0$ и $p_i=E_i$, поэтому
\begin{equation*}
    m_{ij}^2=2E_iE_j(1-\cos\theta_{ij}).
\end{equation*}
В частности,
\begin{align*}
    m_{12}^2&=2\cdot1\cdot1.5(1-\cos6.3^\circ)
    \simeq0.0181\,\text{ГэВ}^2,\\
    m_{12}&\simeq0.135\,\text{ГэВ},\\[0.4em]
    m_{34}^2&=2\cdot1.7\cdot0.5(1-\cos34.6^\circ)
    \simeq0.301\,\text{ГэВ}^2,\\
    m_{34}&\simeq0.548\,\text{ГэВ}.
\end{align*}
Для полноты результаты всех пар:
\begin{center}
    \begin{tabular}{ccccccc}
        \hline
        \rconcept{Пара} & 12 & 13 & 14 & 23 & 24 & 34\\
        \hline
        $m_{ij},\ \text{ГэВ}$ & 0.135 & 0.288 & 1.393 & 0.650 & 1.170 & 0.548\\
        \hline
    \end{tabular}
\end{center}

\answer
Пары $(1,2)$ и $(3,4)$ по инвариантным массам соответствуют распадам
\begin{equation*}
    \pi^0\longrightarrow\gamma_1\gamma_2,
    \qquad
    \eta\longrightarrow\gamma_3\gamma_4.
\end{equation*}
При восстановлении единиц это массы около
$135\,\text{МэВ}/c^2$ и $548\,\text{МэВ}/c^2$.

\task[4]{Потери энергии электрона в остаточном газе}
\label{task:vacuum_losses_sem5}

\condition
В кольцевой трубе длиной $\ell=3\,\text{см}$ находится воздух при давлении
$P=10^{-7}\,\text{Па}$. Электрон имеет энергию порядка ГэВ.
Оценить время, за которое он теряет $1\%$ начальной энергии
из-за тормозного излучения в остаточном газе.

\solution
\subsection*{Плотность газа}

Для воздуха используем
\begin{equation*}
    X_0=36.62\,\text{г}/\text{см}^2,\qquad
    \rho_0\simeq1.2\cdot10^{-3}\,\text{г}/\text{см}^3,
    \qquad P_0\simeq10^5\,\text{Па}.
\end{equation*}
При \rassumption{неизменной температуре} плотность идеального газа
пропорциональна давлению:
\begin{equation*}
    \rho=\rho_0\frac{P}{P_0}
    =1.2\cdot10^{-3}\frac{10^{-7}}{10^5}
    =1.2\cdot10^{-15}\,\text{г}/\text{см}^3.
\end{equation*}

\subsection*{Суммарный путь до потери одного процента энергии}

Для высокоэнергетического электрона средние радиационные потери задаются
\begin{equation*}
    -\frac{dE}{E}=\frac{dX}{X_0},\qquad X=\rho L,
\end{equation*}
где $L$ --- суммарный путь за много оборотов, а не длина трубы.
После интегрирования
\begin{equation*}
    E=E_0e^{-X/X_0},\qquad
    X=X_0\ln\frac{E_0}{E}.
\end{equation*}
Потере $1\%$ энергии соответствует $E/E_0=0.99$, следовательно,
\begin{equation*}
    L_{1\%}=\frac{X_0}{\rho}\ln\frac1{0.99}
    \simeq\frac{36.62}{1.2\cdot10^{-15}}\cdot0.01005
    \simeq3.07\cdot10^{14}\,\text{см}.
\end{equation*}

\subsection*{Время движения по кольцу}

Электрон с энергией порядка ГэВ ультрарелятивистский, поэтому $v\simeq c$.
Один оборот занимает
\begin{equation*}
    t_{\text{об}}\simeq\frac{\ell}{c}
    \simeq10^{-10}\,\text{с}.
\end{equation*}
Число оборотов до указанной потери энергии:
\begin{equation*}
    n\simeq\frac{L_{1\%}}{\ell}\simeq1.02\cdot10^{14}.
\end{equation*}
Таким образом,
\begin{equation*}
    t_{1\%}=nt_{\text{об}}\simeq\frac{L_{1\%}}c
    \simeq1.02\cdot10^4\,\text{с}\simeq2.8\,\text{ч}.
\end{equation*}

\answer
Время потери $1\%$ средней энергии в остаточном газе составляет
\rresult{порядка $10^4\,\text{с}$, или $2.8$ часа}.
При постоянной плотности газа длина кольца сокращается в оценке времени,
но определяет число оборотов.

\end{document}
